% Exercices issus de 7-IntroductionSA2022.docx.
\SAimage{images/image39.png}{70mm}{95mm}

\SAimage{images/image40.jpeg}{158mm}{105mm}

\SAimage{images/image410.png}{22mm}{22mm}
Niveau 1

\textbf{BRAS DE ROBOT PELLENC}

Caractériser la réponse du système non corrigé (précision, rapidité,
dépassement).

Consigne

Réponse (position

du bras)

Position du

bras (

$^\circ$

)

59.5

0

0

10

20

30

40

50

60

temps (s)

1.8

2

1.6

1.4

1.2

1

0.8

0.6

0.4

0.2

Caractériser la réponse du système corrigé (précision, rapidité,
dépassement).

Position du

bras (

$^\circ$

)

52,5

0

0.45

0.5

0

temps (s)

0.4

0.35

0.3

0.25

0.2

0.15

0.1

0.05

10

20

30

40

50

60

\textbf{CORDEUSE DE RAQUETTES}

\SAimage{images/image410.png}{22mm}{22mm}
Niveau 1

Caractériser la réponse du système non corrigé (précision, rapidité,
dépassement).

\SAimage{images/image42.jpeg}{150mm}{75mm}

Consigne

Réponse (tension

de la corde)

Tension de la

corde (N)

133

0

1

2

3

4

5

6

7

8

9

10

0

50

100

150

200

250

300

temps (s)

Caractériser la réponse du système corrigé (précision, rapidité,
dépassement).

Tension de la

corde (N)

271

0

0.9

1

0

50

100

150

200

250

300

temps (s)

0.8

0.7

0.6

0.5

0.4

0.3

0.2

0.1

\SAimage{images/image43.jpeg}{30mm}{42mm}
\textbf{ROBOT D'INSPECTION TUBULAIRE}

\SAimage{images/image410.png}{22mm}{22mm}
Niveau 2

Caractériser les réponses suivantes(rapidité, dépassement).

\SAimage{images/image44.png}{155mm}{92mm}

\begin{quote}
\emph{\textbf{Réponse en déplacement en réponse à un échelon de
pression}}

\SAimage{images/image45.png}{150mm}{110mm}

\emph{\textbf{Réponse de l'ensemble (servovalve + actionneur)}}
\end{quote}

\textbf{20}

\textbf{\ul{Nom~:}}

\begin{enumerate}
\def\labelenumi{\arabic{enumi}.}
\item
  ~Donner les critères permettant d'évaluer un système asservi.
  \emph{\textbf{(2 points)}}
\end{enumerate}

\begin{enumerate}
\def\labelenumi{\arabic{enumi}.}
\item
  ~~Indiquer la différence entre un système bouclé et un système non
  bouclé. \emph{\textbf{(2 points)}}
\item
  Donner les deux types de systèmes asservis. \emph{\textbf{(1 points)}}
\end{enumerate}

\SAimage{images/image46.png}{155mm}{95mm}

\textbf{1,75}

\begin{enumerate}
\def\labelenumi{\arabic{enumi}.}
\setcounter{enumi}{1}
\item
  Donner le temps de réponse à 5 \% t\textsubscript{r5\%} et l'écart
  statique. \emph{\textbf{(2 points)}}
\end{enumerate}

\SAimage{images/image47.png}{145mm}{80mm}

\begin{enumerate}
\def\labelenumi{\arabic{enumi}.}
\setcounter{enumi}{2}
\item
  Donner le temps de réponse à 5 \%, le dépassement et l'écart statique.
  \emph{\textbf{(3 points)}}
\end{enumerate}

\SAimage{images/image48.png}{55mm}{38mm}

\SAimage{images/image410.png}{22mm}{22mm}
Niveau 2

\textbf{RÉGULATION DE NIVEAU D'EAU D'UN RÉSERVOIR D'INCENDIE}

\emph{(\ul{Source :} BTS IRIS 2008)}

\textbf{\ul{Objectifs de l'étude~:}}

\begin{itemize}
\item
  \textbf{Modéliser} le réservoir d'incendie à partir du modèle de
  connaissance ;
\item
  \textbf{Établir} le schéma-bloc du système~;
\end{itemize}

Pour pouvoir faire face en cas d'incendie, de nombreuses communes
disposent d'un réservoir d'eau. Le réservoir étudié est alimenté en eau
par une pompe entraînée par une machine asynchrone alimentée par un
variateur de fréquence~: ensemble redresseur + onduleur. La hauteur
d'eau est mesurée par un capteur de pression placé au fond du réservoir.
On peut schématiser le fonctionnement du système par la figure
suivante~:

\SAimage{images/image49.jpeg}{150mm}{82mm}

Le réservoir est un cylindre de section constante S= 3.5
m\textsuperscript{2}. On considère que le débit de sortie
q\textsubscript{S} (en m\textsuperscript{3}.s\textsuperscript{-1}) est
proportionnel à la hauteur d'eau h (en m).

On désire que le réservoir se remplisse à 95\% de la hauteur voulue en
15mn, avec une précision en régime permanent de 10 \%.

\textbf{Modélisation du réservoir}

On veut vérifier les performances par une étude comportementale, en
faisant à l'instant t = 0, varier le débit de 0 à
0.01m\textsuperscript{3}.s\textsuperscript{-1}. On enregistre
l'évolution de la hauteur h de l'eau en fonction du temps comme le
montre la figure suivante~:

\SAimage{images/image50.png}{150mm}{78mm}

\begin{enumerate}
\def\labelenumi{\arabic{enumi}.}
\item
  \textbf{Indiquer} la valeur \emph{\textbf{h\textsubscript{$\infty$}}} de la
  hauteur d'eau en régime permanent .
\item
  \textbf{Caractériser} cette réponse (rapidité, précision, stabilité,
  dépassement).
\item
  \textbf{Donner} le temps nécessaire pour que le réservoir se remplisse
  à 95\% de la hauteur voulue. Comparer au cahier des charges.

  \textbf{Régulation du niveau d'eau du réservoir}
\end{enumerate}

On désire que le réservoir se remplisse à 95\% de la hauteur voulue en
15mn, avec une précision en régime permanent de 10 \%. Afin d'améliorer
les performances du système, on insère dans l'asservissement un
correcteur proportionnel de transmittance \emph{\textbf{C(p)=
C\textsubscript{0}}}.

On souhaite mettre le système sous forme de schéma-blocs. Les modèles
des différents éléments du système sont les suivants~:

\begin{itemize}
\item
  Le capteur de pression mesurant la hauteur d'eau délivre une tension
  \emph{\textbf{V\textsubscript{h}(t) = $\alpha$ h(t)}} .
\end{itemize}

\begin{itemize}
\item
  La salle de commande élabore la tension de commande \textbf{V(t)} du
  variateur à partir de la différence entre la tension de consigne
  \emph{\textbf{V\textsubscript{c}(t)}} et la mesure provenant du
  capteur \emph{\textbf{V\textsubscript{h}(t)}}. \textbf{V(t)=
  C\textsubscript{0} (\emph{V\textsubscript{c}(t)-
  V\textsubscript{h}(t)}}).
\end{itemize}

\begin{itemize}
\item
  La relation entre la hauteur de consigne
  \emph{\textbf{h\textsubscript{c}(t)}} et la tension de consigne est
  \emph{\textbf{V\textsubscript{c}(t) = $\alpha$ h\textsubscript{c}(t)}}
\end{itemize}

\begin{itemize}
\item
  On modélise le fonctionnement de l'ensemble «~moteur+variateur+pompe~»
  par \emph{\textbf{Q\textsubscript{E}(t) = A .V(t), où
  Q\textsubscript{E}(t)}} représente le débit fourni par la pompe et
  \emph{\textbf{V(t)}} la tension de commande.
\end{itemize}

\begin{enumerate}
\def\labelenumi{\arabic{enumi}.}
\setcounter{enumi}{3}
\item
  \textbf{Donner} le schéma bloc de l'asservissement, ayant pour entrée
  \emph{\textbf{H\textsubscript{c}(p}}) (transformée de Laplace de la
  hauteur de consigne h\textsubscript{c}), et pour sortie
  \emph{\textbf{H(p).}}
\end{enumerate}

\emph{\hfill\break
}

\textbf{FAUTEUIL ELECTRIQUE}

\SAimage{images/image410.png}{22mm}{22mm}
Niveau 2

\SAimage{images/image51.pdf}{155mm}{95mm}

Le système étudié est un fauteuil électrique capable de franchir
obstacles et marches.

\textbf{Structure de régulation de l'inclinaison du siège}

La capacité de monter les escaliers nécessite, pour le confort et la
sécurité du passager, de maintenir l'inclinaison du siège par rapport à
la verticale constante pendant la montée ou la descente, en conservant
la possibilité pour celui-ci de régler l'inclinaison de son assise.

Le choix fait par le concepteur est de réaliser un asservissement de la
position angulaire du siège à l'aide d'un vérin électrique. Le but de
cette partie est de déterminer la structure de cet asservissement en
partie décrite par le diagramme IBD limité à la chaîne asservie (voir
annexes).

\ul{La partie commande} de cet asservissement comporte :

\begin{itemize}
\item
  Une interface de communication avec l'usager (bloc « IHM »)
\end{itemize}

Siège

B

Ecrou

Vis

Corps

A

O

Châssis

\begin{itemize}
\item
  Une carte de commande dont le cœur est « l'unité de traitement » qui
  ne reçoit et ne délivre que des signaux numériques
\item
  Un accéléromètre délivrant un signal analogique image de l'inclinaison
  du siège par rapport à la verticale
\end{itemize}

\ul{La partie opérative} (voir schéma cinématique mécanique ci-contre)
de cet asservissement comporte :

\begin{itemize}
\item
  Un pont en H (hacheur quatre quadrants) qui à partir du signal
  analogique reçu de la carte de commande et de l'énergie fournie par un
  régulateur de tension alimente le motoréducteur du vérin
\item
  Un motoréducteur (moteur à courant continu et réducteur à engrenages)
  qui fait tourner la vis du vérin
\item
  Un système vis écrou permettant de transformer la rotation de la vis
  en un allongement du vérin (AB)
\item
  Un système triangulaire (ABO) modifiant l'inclinaison du siège
\end{itemize}

L'énergie électrique est fournie par deux batteries de 12V 60Ah dont la
tension diminue de manière non négligeable au fur et à mesure que la
batterie se décharge

\begin{enumerate}
\def\labelenumi{\arabic{enumi}.}
\item
  Indiquer ci-dessous à quels éléments correspondent les blocs du
  diagramme IBD.
\end{enumerate}

\begin{longtable}[]{@{}
  >{\raggedright\arraybackslash}p{(\columnwidth - 8\tabcolsep) * \real{0.0968}}
  >{\raggedright\arraybackslash}p{(\columnwidth - 8\tabcolsep) * \real{0.2258}}
  >{\raggedright\arraybackslash}p{(\columnwidth - 8\tabcolsep) * \real{0.2258}}
  >{\raggedright\arraybackslash}p{(\columnwidth - 8\tabcolsep) * \real{0.2258}}
  >{\raggedright\arraybackslash}p{(\columnwidth - 8\tabcolsep) * \real{0.2258}}@{}}
\toprule\noalign{}
\begin{minipage}[b]{\linewidth}\raggedright
Bloc
\end{minipage} & \begin{minipage}[b]{\linewidth}\raggedright
Modulateur
\end{minipage} & \begin{minipage}[b]{\linewidth}\raggedright
Convertisseur
\end{minipage} & \begin{minipage}[b]{\linewidth}\raggedright
Transmetteur
\end{minipage} & \begin{minipage}[b]{\linewidth}\raggedright
Effecteur
\end{minipage} \\
\midrule\noalign{}
\endhead
\bottomrule\noalign{}
\endlastfoot
Elément & & & & \\
\end{longtable}

\emph{On donne dans le tableau ci-dessous les noms des différents flux
d'information et d'énergie du diagramme IBD de l'annexe}

On suppose que la carte de commande est constituée d'un comparateur qui
compare la tension u\textsubscript{C}(t) image de la consigne de
position du siège et la tension u\textsubscript{A}(t) image de la
position angulaire du siège issue de l'accéléromètre.

L'écart entre ces deux positions étant ensuite corrigé par un correcteur
qui délivre au pont en H une tension u\textsubscript{H}(t). Lequel pont
en H alimente le moteur avec la tension modulée u\textsubscript{M}(t).

On adopte les notations des fonctions temporelles du système suivantes :

\begin{longtable}[]{@{}
  >{\raggedright\arraybackslash}p{(\columnwidth - 6\tabcolsep) * \real{0.4220}}
  >{\raggedright\arraybackslash}p{(\columnwidth - 6\tabcolsep) * \real{0.0762}}
  >{\raggedright\arraybackslash}p{(\columnwidth - 6\tabcolsep) * \real{0.4219}}
  >{\raggedright\arraybackslash}p{(\columnwidth - 6\tabcolsep) * \real{0.0799}}@{}}
\toprule\noalign{}
\begin{minipage}[b]{\linewidth}\raggedright
Tension image de la consigne de position
\end{minipage} & \begin{minipage}[b]{\linewidth}\raggedright
u\textsubscript{C}(t)
\end{minipage} & \begin{minipage}[b]{\linewidth}\raggedright
Tension image de la réponse
\end{minipage} & \begin{minipage}[b]{\linewidth}\raggedright
u\textsubscript{A}(t)
\end{minipage} \\
\midrule\noalign{}
\endhead
\bottomrule\noalign{}
\endlastfoot
Tension d'alimentation du moteur & u\textsubscript{M}(t) & Position
angulaire du siège / à la verticale & $\theta$(t) \\
Tension de commande du pont en H & u\textsubscript{H}(t) & Position
angulaire du siège / au châssis & $\alpha$(t) \\
Vitesse de rotation de la vis & $\omega$\textsubscript{V}(t) & Position
angulaire du châssis / à la verticale & $\beta$(t) \\
Ecart à l'entrée du correcteur & $\varepsilon$(t) & Longueur du vérin & $\lambda$(t) \\
\end{longtable}

\emph{Les noms des différents éléments du système sont : Motoréducteur,
Correcteur, Pont en H, Accéléromètre, Système vis-écrou, Système
triangulaire}

\begin{enumerate}
\def\labelenumi{\arabic{enumi}.}
\setcounter{enumi}{1}
\item
  Placer les noms des différents blocs et les fonctions temporelles sur
  le schéma-bloc ci-dessous.
\end{enumerate}

\SASchemaFauteuil

\textbf{\ul{Modélisation du vérin}}

\emph{Le} vérin électrique comprend le motoréducteur et le système vis
écrou. Il est constitué de :

\begin{itemize}
\item
  Un moteur à courant continu à aimants permanents pour lequel :
\end{itemize}

\begin{itemize}
\item
  On note u\textsubscript{M}(t) et i(t) les tensions et courant aux
  bornes du moteur
\item
  On note e(t) la force contre électromotrice du moteur
\item
  On note $\omega$\textsubscript{M}(t) et c\textsubscript{M}(t) la vitesse de
  rotation de l'arbre moteur et le couple sur l'arbre moteur
\item
  On néglige l'inductance de l'induit L mais pas sa résistance R
\item
  L'inertie de l'ensemble des pièces en mouvement ramené sur l'arbre
  moteur est J
\item
  Le coefficient de frottement visqueux des pièces en mouvement ramené
  sur l'arbre moteur est f
\item
  Les constantes mécanique et électrique sont K\textsubscript{M} et
  K\textsubscript{E}
\end{itemize}

\begin{itemize}
\item
  Un réducteur à engrenages qui transmet la rotation de l'arbre moteur à
  la vis. On note r le rapport de transmission de ce
  réducteur,$\omega$\textsubscript{M}(t) et $\omega$\textsubscript{V}(t)les vitesses
  de rotation de l'arbre moteur et de la vis.
\item
  Un système vis écrou qui transforme la vitesse de rotation de la vis
  $\omega$\textsubscript{V}(t) en allongement du vérin $\lambda$(t).On note « pas » le
  pas de la vis (distance par courue par l'écrou pour un tour de la vis)
\end{itemize}

\begin{enumerate}
\def\labelenumi{\arabic{enumi}.}
\setcounter{enumi}{2}
\item
  Donner la relation entre $\omega$\textsubscript{V}(t) et
  $\omega$\textsubscript{M}(t) d'une part, celle entre $\omega$\textsubscript{V}(t) et
  $\lambda$(t) d'autre part.
\end{enumerate}

\textbf{\ul{Analyse des performances de l'asservissement en fonction du
correcteur}}

Pour valider le comportement de cet asservissement le concepteur a
installé un correcteur qui possède deux paramètres de réglage~: un gain
K et un coefficient a. Il a réalisé quatre essais avec des valeurs
différentes de a et K avec le protocole suivant :

\begin{itemize}
\item
  Le fauteuil est placé sur un sol horizontal ($\beta$(t) = 0)
\item
  Un échelon de consigne de $\theta$\textsubscript{C} = --5$^\circ$ est appliqué au
  système
\item
  L'évolution temporelle de la réponse $\theta$(t) est relevée
\end{itemize}

On donne ci-dessous les différentes réponses temporelles pour ces quatre
essais. On a en ordonnée la réponse $\theta$\textsubscript{C} en degré et en
abscisse le temps t en seconde.

\begin{enumerate}
\def\labelenumi{\arabic{enumi}.}
\setcounter{enumi}{3}
\item
  Calculer et donner ci-dessous, pour chacun de ces essais, l'erreur
  statique relative $\varepsilon$\textsubscript{S}, le temps de réponse à 5\%
  t\textsubscript{r5\%}, et le taux de dépassement D\textsubscript{1\%}
  si il existe.
\end{enumerate}

\begin{longtable}[]{@{}
  >{\raggedright\arraybackslash}p{(\columnwidth - 2\tabcolsep) * \real{0.5083}}
  >{\raggedright\arraybackslash}p{(\columnwidth - 2\tabcolsep) * \real{0.4917}}@{}}
\toprule\noalign{}
\begin{minipage}[b]{\linewidth}\raggedright
\SAimage{images/image55.png}{75mm}{65mm}
\end{minipage} & \begin{minipage}[b]{\linewidth}\raggedright
\SAimage{images/image56.png}{75mm}{65mm}
\end{minipage} \\
\midrule\noalign{}
\endhead
\bottomrule\noalign{}
\endlastfoot
\SAimage{images/image57.png}{75mm}{65mm} &
\SAimage{images/image58.png}{75mm}{65mm} \\
\end{longtable}

\begin{enumerate}
\def\labelenumi{\arabic{enumi}.}
\setcounter{enumi}{4}
\item
  En vous aidant du diagramme d'exigence (voir annexes), dire quels sont
  les cas qui respectent le cahier des charges relatif à l'inclinaison
  de l'assise. Argumenter la réponse.
\end{enumerate}

\SAimage{images/image59.png}{150mm}{100mm}

Diagramme de blocs interne limité à la chaîne asservie d'inclinaison du
siège

\SAimage{images/image60.png}{150mm}{165mm}

Diagramme d'exigence de premier niveau du fauteuil
